\psubsection{fairness}{0.9}{I2}{legal,core}{Fairness as a Policy Choice}
% VI-B, contribution I2 (brief private/VERA_ICSE2027_writing_brief.md (not versioned), section VI-B).
% Scenarios (decision D4, 30 Sep):
%   full       P1, P2 with Proposition 1, P3 (tabular results), P4 (N2), P6, P7;
%              Fig. 2 and Tab. 3 with eight rows.
%   reduced    P3 on the synthetic corpus only, point estimates; P4 carries the
%              criterion analysis; Tab. 3 has five rows; brief budget 0.8 p.
%   withdrawn  P3 and Fig. 2 removed; P4 (N2) carries I2 alone through Tab. 3
%              (language-model rows); brief budget 0.6 p. If N2 does not run (D7),
%              also drop P4 and Tab. 3 and state I2 as a design contribution.
% Brief P5 (E7 reweighting of the R10 benchmarks) is left out: the R10 instruments
% are not valid measurements as they stand (brief 0.2). One sentence goes to IX
% unless D7 repairs them.
% Proposition 1 is not in the APSEC paper, which leaves weight dependence out of
% scope. Never present the questionnaire's fairness flag as an output of it.

\armnotwithdrawn{%
\plannedfloat{figure}{fig:criterion-bands}{\columnwidth}{4.5cm}%
  {Reachable fairness range over the admissible criteria \armswitch{for each tabular model--dataset pair and each language-model deployment}{for the synthetic corpus and each language-model deployment}{}, against the review bands. Ticks mark the score of each criterion; a cross marks a criterion-dependent verdict (\Cref{prop:criterion}).}%
  {Generated by N3 (gap G12): generalise gen\_sensitivity\_panel.py and the vera-foundry E7 band figure from benchmark weights to criteria, output figures/fairness\_criterion\_bands.pdf. One bar per row of Tab.~3 from the lowest to the highest criterion score, over the red, orange and green bands, one tick per criterion, a cross where the two ends fall in different bands; tabular models and language models as two facets.}%
}

\generatedtable{table}{fairness_criteria}{tab:fairness-criteria}{\columnwidth}{3.6cm}%
  {Fairness score per admissible criterion, and criterion robustness}%
  {Generated by N3 from the N1 tabular runs and the N2 language-model runs into tables/fairness\_criteria.tex, label tab:fairness-criteria. Rows: synthetic corpus (planted truth, reference model), Adult, South German Credit, language-model deployments. Columns: DI, EO, EOd, gECE, CF (if kept, D3), range, robust, proxy (if kept, D3).}

% P1 -- the impossibility results, and why a review must name its criterion.
Outside degenerate cases, the standard group-fairness criteria conflict.
For a risk score, calibration within groups and balance of the average score across groups within each true class can hold together only under perfect prediction or equal base rates~\cite{kleinberg2017inherent}.
For an imperfect classifier and groups with different base rates, equal positive predictive value cannot coexist with equal false positive and equal false negative rates~\cite{chouldechova2017fair}.
Calibration is compatible with at most one error-rate constraint~\cite{pleiss2017fairness}.
Choosing among them is normative~\cite{barocas2023fairness}, so a review must name its criterion.

% P2 -- the criterion as a weighting, and Proposition 1.
% IV (04-vera) already states that each criterion is a benchmark under R11 and
% that choosing one is a one-hot weighting; do not repeat it here.
Each admissible criterion is a benchmark under \req{11} (\Cref{sec:vera}) whose score lies in $[0,1]$, higher meaning fairer, such as one minus the true-positive-rate gap for equal opportunity~\cite{hardt2016equality} or the ratio of the lower to the higher rate of favourable decisions~\cite{feldman2015certifying} \decision{D3}{core}{22 Sep}{Keep or cut the attribute-flip (counterfactual) and proxy checks: they decide the CF and Proxy columns of Tab. 3.}.

\begin{proposition}[Criterion robustness]\label{prop:criterion}
Let $\mathcal{C}$ be a finite non-empty set of admissible criteria and $m_c\in[0,1]$ the score of criterion $c\in\mathcal{C}$ for one model and dataset, higher meaning fairer.
A policy is a weighting $w\in\mathcal{W}=\{w\in\mathbb{R}_{\ge 0}^{\mathcal{C}}:\sum_{c}w_c=1\}$ with score $s(w)=\sum_{c}w_c\,m_c$; adopting criterion $c$ alone is the one-hot weighting $e_c$, any other $w$ a mixture.
Let $\beta:[0,1]\to\{\mathrm{red}<\mathrm{orange}<\mathrm{green}\}$ be non-decreasing.
Then $\beta(\min_c m_c)\le\beta(s(w))\le\beta(\max_c m_c)$ for every $w\in\mathcal{W}$, both bounds are attained, and the verdict $\beta(s(w))$ is the same for every criterion and every mixture, that is, \emph{criterion-robust}, if and only if $\beta(\min_c m_c)=\beta(\max_c m_c)$.
\end{proposition}
\begin{IEEEproof}
$s(w)$ is a convex combination of the $m_c$, so $\min_c m_c\le s(w)\le\max_c m_c$, and a non-decreasing $\beta$ keeps this order.
The one-hot weightings on a minimising and on a maximising criterion attain the two bounds, so all verdicts coincide exactly when the two bounding bands do.
\end{IEEEproof}

% P2, continued -- common scale, per-criterion thresholds, interval variant.
% The interval variant needs row-level bootstrap intervals (N1, gap G4, 30 Sep).
% It is dropped in the reduced arm (point estimates); if G4 does not land, delete
% it in every arm and state point estimates as a threat in IX.
VERA applies the same two band thresholds, set in the deployment environment, to every requirement, which is the common scale assumed here; its sensitivity scripts already run this band test over benchmark weights.
With a threshold $\tau_c$ per criterion instead \decision{D5}{legal}{30 Sep}{Admissible criteria for credit decisions; one common scale or a threshold per criterion; default criterion; quotable rationale.}, the verdict survives every single-criterion choice if and only if the decisions $m_c\ge\tau_c$ agree over $\mathcal{C}$, which needs neither convexity nor a common scale.
\ifarmreduced\else
Taking the lowest lower and the highest upper bootstrap confidence bound~\cite{efron1993} instead of $\min_c m_c$ and $\max_c m_c$ makes the test conservative.
\fi

% P3 -- tabular results (N1, N3). Removed when the ML arm is withdrawn.
\ifarmfull
\Cref{tab:fairness-criteria} gives continuous scores and \Cref{fig:criterion-bands} their ranges.
On the synthetic corpus, the reference model's disparate-impact score is \tbd{N1}{core}{04 Oct}{disparate-impact score of the reference model on the synthetic corpus} against a planted \tbd{N1}{core}{04 Oct}{disparate-impact ratio planted by the synthetic generator (G9 manifest)}; \tbd{N3}{core}{09 Oct}{number of criterion-dependent tabular verdicts} of \tbd{N3}{core}{09 Oct}{number of tabular model--dataset pairs} tabular verdicts are criterion-dependent \decision{D6}{legal}{22 Sep}{Protected attribute(s), sex or age, and South German Credit instead of the UCI coding.}.
\parplan[12]{Adult and South German Credit are set against credit-scoring fairness studies.}{N3 on Adult (becker1996adult) and South German Credit (hofmann1994statlog with gromping2019south; D6). Continuous values only, never a thresholded 0 or 1. Contrast with kozodoi2022fairness (criteria compared on credit data), hurlin2026fairness (fairness tests for scoring models) and biswas2020machine (metric disagreement on real models). Point estimates unless the row-level bootstrap (N1, gap G4) lands by 30 Sep; otherwise a threat in IX.}
\fi
\ifarmreduced
\Cref{tab:fairness-criteria} gives continuous point estimates and \Cref{fig:criterion-bands} their ranges.
The tabular side covers the synthetic corpus only: the reference model's disparate-impact score is \tbd{N1}{core}{04 Oct}{disparate-impact score of the reference model on the synthetic corpus} against a planted \tbd{N1}{core}{04 Oct}{disparate-impact ratio planted by the synthetic generator (G9 manifest)}, and its verdict is \tbd{N3}{core}{09 Oct}{criterion-robust or criterion-dependent verdict on the synthetic corpus} \decision{D6}{legal}{22 Sep}{Protected attribute(s), sex or age, of the synthetic corpus.}.
\fi

% P4 -- language models classifying credit records (N2).
\ifarmwithdrawn
\Cref{tab:fairness-criteria} applies the test to language models that classify credit records rendered as text, whose decisions feed the same metrics module \decision{D7}{core}{22 Sep}{Language-model fairness instrument: N2 record classification, or keep the R10 and R11 probes with caveats (then I2 is a design contribution).}: \tbd{N2}{foundry}{04 Oct}{number of criterion-dependent verdicts among the deployments run under N2} of \tbd{N2}{foundry}{04 Oct}{number of deployments run under N2} deployments have a criterion-dependent verdict.
\else
When a language model classifies credit records rendered as text, the same module and test apply \decision{D7}{core}{22 Sep}{Language-model fairness instrument: N2 record classification, or keep the R10 and R11 probes with caveats.}: \tbd{N2}{foundry}{04 Oct}{number of criterion-dependent verdicts among the deployments run under N2} of \tbd{N2}{foundry}{04 Oct}{number of deployments run under N2} deployments have a criterion-dependent verdict.
\fi
% II (02-background) and IX (09-threats) cite the benchmark-validity critique; not repeated here.
The engine's language-model instruments for \req{10} and \req{11} are not used as fairness evidence (\Cref{sec:threats}): the \req{11} probe renders one prompt template and its score reduces to the share of positive answers.
\ifarmreduced
Record classification has outcome labels, so every error-rate criterion is computable for each deployment; the criteria that separate the bands are \tbd{N3}{core}{09 Oct}{criteria that separate the bands, per deployment run under N2}.
\fi
\ifarmwithdrawn
The criteria that separate the bands are \tbd{N3}{core}{09 Oct}{criteria that separate the bands, per deployment run under N2}.
\fi

% P6 -- the second line owns the criterion (N10, D5, G5).
The second line of defence owns this choice.
EU non-discrimination law resists reduction to one statistical test~\cite{wachter2021fairness}, and the revised Consumer Credit Directive, applicable from 20~November 2026, prohibits discrimination in access to credit without naming a metric~\cite{euccd2023}.
For credit decisions, \legalrisk{} admitted \tbd{N10}{legal}{30 Sep}{admissible criteria for credit decisions}, chose \tbd{N10}{legal}{30 Sep}{common scale or per-criterion thresholds, and the default criterion} and justified the choice as \tbd{N10}{legal}{30 Sep}{quotable rationale for the criterion policy}.
The tabular specification records every admissible criterion under \req{11} in its catalog, the non-default ones with zero weight \tbd{N1}{core}{25 Sep}{tabular catalog listing every admissible criterion under R11 (gap G11)}, so the content digest (SHA-256) pinned into each run covers the criterion policy; the band thresholds are environment settings outside it \decision{D29}{core}{25 Sep}{G5: put thresholds and registry into the catalog and its digest, or keep describing the thresholds as environment settings.}.

% P7 -- what a review reports.
A review thus reports whether each fairness verdict is criterion-robust and, if not, which criteria move it across a band, not a single fairness score.
